% 03_methods.tex
\section{Methods}
\label{sec:methods}

\subsection{Data}
\label{sec:methods_data}

We use the complete \textit{H.~sapiens} proteome from AlphaFold Protein
Structure Database version~4~\cite{varadi2022alphafold}, distributed as a
single tar archive \texttt{UP000005640\_9606\_HUMAN\_v4.tar} containing
23{,}391 gzip-compressed PDB files.
Each file corresponds to one AlphaFold fragment (F1, F2, \ldots{} for
multi-fragment proteins) and is treated as a distinct structure throughout,
consistent with the paper's fragment-level semantics.
Per-residue pLDDT values are read from the B-factor column of PDB ATOM
records using the gemmi library~\cite{wojdyr2022gemmi}.

\paragraph{Fragment length filter.}
Following the paper's Fig.~7--8 analysis, we restrict all quantitative
comparisons to first-fragment (F1) entries with $n \geq 200$ residues
(16{,}787 of the 23{,}391 fragments).
Shorter fragments tend to have near-zero $H_p$ regardless of sequence
complexity, as a small persistence spectrum trivially has few distinct
values.

\paragraph{pLDDT discretisation.}
AlphaFold-DB v4 stores pLDDT as floating-point values.
A 316-residue protein such as P15121 can exhibit 108 distinct raw pLDDT
values, resulting in a persistence spectrum with many singleton
persistence values and artificially inflated $H_p$.
To match the paper's regime — which appears to use effectively
integer-valued pLDDT distributions ($\lesssim\! 20$ distinct levels per
protein) — we apply the rounding function $\ell_i = \lfloor s_i +
0.5\rfloor$ before the filtration.
This reduces P15121 to 13 distinct levels and brings $H_p$ within
${\approx}0.04$ of the paper's target.

\subsection{Path-graph filtration}
\label{sec:methods_filt}

We implement the filtration from scratch using Union-Find with path-halving
and the Elder Rule.
A key design choice is \emph{batch insertion}: all residues with the same
discretised pLDDT level $L$ are inserted as a group before any
cross-level edges are activated.
This ensures that intra-batch edges produce accretion pairs ($b=d=L$),
while inter-batch edges produce positive-persistence pairs, faithfully
capturing the multi-scale pLDDT structure.

\medskip
\noindent\rule{\linewidth}{0.6pt}
\smallskip
\noindent\textbf{Algorithm~1:} Batch-insertion pLDDT path-graph filtration.
\smallskip
\noindent\rule{\linewidth}{0.4pt}
\smallskip
\noindent\begin{minipage}{\linewidth}\small
\noindent\textbf{Input:} pLDDT array $(s_1,\ldots,s_n)$;
  discretiser $\delta$ (default: $\delta(x) = \lfloor x + 0.5\rfloor$)\\
\noindent\textbf{Output:} persistence diagram $\mathcal{P}$;
  $N_{cc}$-curve (pairs of level, count)\\[4pt]
\noindent Compute $\ell_i \leftarrow \delta(s_i)$; form sorted distinct
  levels $L_1 > L_2 > \cdots > L_k$\\[2pt]
\noindent\textbf{for} each level $L \in (L_1,\ldots,L_k)$ \textbf{do}\\
\hspace*{1.2em}$B \leftarrow \{i : \ell_i = L\}$\quad
  \textit{// all residues at this level}\\
\hspace*{1.2em}\textbf{for} $i \in B$ (asc.\ index):
  insert singleton $i$;\ $\mathrm{birth}[i]\!\leftarrow\! L$;\
  $N_{cc}\!\mathrel{+}=\!1$\\
\hspace*{1.2em}\textbf{for} $i \in B$ (asc.\ index):\\
\hspace*{2.4em}\textbf{if} $i > 1$ \textbf{and} $i{-}1$ is inserted:
  \textsc{Union}$(i,\;i{-}1,\;L)$\\
\hspace*{2.4em}\textbf{if} $i < n$ \textbf{and} $i{+}1$ is inserted
  \textbf{and} $\ell_{i+1} > L$: \textsc{Union}$(i,\;i{+}1,\;L)$\\
\hspace*{1.2em}Append $(L,\,N_{cc})$ to $N_{cc}$-curve for each
  $i \in B$\\[4pt]
\noindent\textbf{procedure} \textsc{Union}$(a,\,b,\,d)$:\quad
  \textit{// $d$ = death pLDDT}\\
\hspace*{1.2em}$r_a \leftarrow \textsc{Find}(a)$;\quad
  $r_b \leftarrow \textsc{Find}(b)$\\
\hspace*{1.2em}\textbf{if} $r_a = r_b$: append $(d,\,d)$ to $\mathcal{P}$;
  \textbf{return}\quad\textit{// accretion}\\
\hspace*{1.2em}\textit{survivor} $\leftarrow$ root with higher
  $\mathrm{birth}[\cdot]$; tie-break: lower index\quad
  \textit{// Elder Rule}\\
\hspace*{1.2em}$\mathrm{parent}[\textit{loser}]\leftarrow\textit{survivor}$;\quad
  $N_{cc}\!\mathrel{-}=\!1$\\
\hspace*{1.2em}Append $(\mathrm{birth}[\textit{loser}],\;d)$ to $\mathcal{P}$
\end{minipage}
\smallskip
\noindent\rule{\linewidth}{0.6pt}
\medskip

\paragraph{Edge-processing invariant.}
Each path-graph edge is processed exactly once.
For batch vertex $i$, the left edge $(i{-}1, i)$ is checked if $i{-}1$ is
already inserted (covering both within-batch and cross-level left neighbours).
The right edge $(i, i{+}1)$ is only activated when $i{+}1$ already belongs
to a \emph{higher} level ($\ell_{i+1} > L$); within-batch right neighbours
will encounter their own left-check on their turn, preventing double processing.

\paragraph{Correctness on P15121.}
For the 316-residue ordered protein P15121 the algorithm produces $H_p = 0.08$
(paper target: ${\approx}0.04$); the residual offset of $0.04$ is consistent
with the number of remaining float-precision levels after rounding.

\subsection{Five statistics}
\label{sec:methods_stats}

All five statistics are computed from $\mathcal{P}$ and the $N_{cc}$-curve
returned by Algorithm~1.

\paragraph{$f^+_{cp}$}
is computed as $\sum_i\mathbf{1}[b_i > d_i]\,/\,(n{-}1)$ via a single
NumPy boolean comparison on the two columns of $\mathcal{P}$~\cite{harris2020numpy}.

\paragraph{$\bar{p}$}
is the mean of the column-difference $\mathcal{P}[:,0] - \mathcal{P}[:,1]$.

\paragraph{$H_p$}
is computed from Eq.~\eqref{eq:Hp} using \texttt{numpy.unique(...,
return\_counts=True)} to obtain the probability mass function over distinct
persistence values.
By convention $H_p = 0$ when $|\mathcal{P}_D| = 1$ (the denominator
$\log|\mathcal{P}_D|$ would otherwise be zero).

\paragraph{$N_{cc}^{max}$}
is the maximum of the second column of the $N_{cc}$-curve.

\paragraph{$\mathrm{PLM}(t_p)$}
is computed as described in Section~\ref{sec:methods_plm} below.

\subsection{PLM via GUDHI}
\label{sec:methods_plm}

Persistent local maxima of the scalar function $N_{cc}(u)$ are identified
via the Morse-Smale chain-complex simplification described
in~\cite{cazals2025alphafold}.
We implement this using the GUDHI library's
\texttt{CubicalComplex}~\cite{gudhi2015}:

\begin{enumerate}
  \item Extract the $N_{cc}$ column of the $N_{cc}$-curve and negate it:
        $\widetilde{N} = -N_{cc}$.
        Negation maps local maxima of $N_{cc}$ to local \emph{minima} of
        $\widetilde{N}$, which are the 0-dimensional features born first in
        the sub-level-set filtration.

  \item Construct a \texttt{CubicalComplex} with
        \texttt{top\_dimensional\_cells} set to $\widetilde{N}$,
        then call \texttt{persistence()}.

  \item For each 0-dimensional pair $(b_j, d_j)$ in the persistence
        diagram of $\widetilde{N}$, the persistence is
        $d_j - b_j = N_{cc}[\text{peak}] - N_{cc}[\text{valley}]$,
        measured in connected-component count units.

  \item Count pairs with $d_j - b_j \geq t_\nu = n \cdot t_p$, where
        $t_p = 0.025$ by default ($t_p \in \{0.020, 0.025, 0.030\}$ for
        the ablation study in Section~\ref{sec:results}).
\end{enumerate}

For the $t_p$-ablation pass, we build the \texttt{CubicalComplex} once and
apply all three thresholds to the same persistence list (\texttt{plm\_multi}),
reducing the GUDHI overhead by $3\times$.

\subsection{Arity extension}
\label{sec:methods_arity}

$\mathrm{C}_\alpha$ coordinates are extracted from the gemmi structure
objects using the same tarball-streaming pipeline.
We build a \texttt{scipy.spatial.cKDTree}~\cite{virtanen2020scipy} on the
$n\times 3$ coordinate array and query it with
\texttt{query\_ball\_point(ca\_coords, r=10.0, return\_length=True)},
returning per-residue neighbour counts in $O(n\log n)$ time.
Subtracting 1 removes the self-count.
The arity signature $(a_{25}, a_{75})$ is computed as integer-truncated
25th and 75th percentiles via \texttt{numpy.percentile}.

\subsection{Proteome-scale pipeline}
\label{sec:methods_pipeline}

The tar archive is streamed in a single sequential pass over its 23{,}391
members using Python's \texttt{tarfile} in streaming mode
(\texttt{mode="r|"}), avoiding the ${\approx}66$\,GB memory cost of random
access.
Each \texttt{.pdb.gz} member is decompressed in memory via
\texttt{gzip.decompress} and dispatched to a pool of six worker processes
(\texttt{multiprocessing.Pool} with \texttt{fork} start method).
Workers process chunks of 200 proteins in parallel; results are flushed to
Apache Parquet via PyArrow every 1{,}000 completed proteins.
Two separate passes produce:
\begin{itemize}
  \item \texttt{proteome\_full.parquet} --- $f^+_{cp}$, $\bar{p}$, $H_p$,
        $N_{cc}^{max}$, $\mathrm{PLM}(0.025)$ for all 23{,}391 fragments;
  \item \texttt{proteome\_full\_arity.parquet} --- $(a_{25}, a_{75})$ for
        all 23{,}391 fragments.
\end{itemize}
The two tables are merged on \texttt{(uniprot\_id, filename)} for the
cross-correlation analysis of Section~\ref{sec:results}.
